% SOURCE AUTHORITY:
% expose_VII_body.tex lines 1810--1977; scan indices 223--228; source folios 212--217.
% Direct scan comparison: physical PDF pages 224--229 at 600 dpi.
% Transparent normalization: typewritten letter-o subscripts are rendered as zero.
% The headless diagonal in (3.8.2) is reproduced as printed, while the prose fixes
% its mathematical role as the trivializing homomorphism.  The derived-category
% arrow in (3.8.10) follows the French authority, from P derived-tensor Q to
% L_\bullet tensor M_\bullet; the inherited abridgment reversed it.  Hard stop at
% the end of the Expose VII bibliography.

\subsection*{3.8. Biextensions of $(P,Q)$ trivialized over extensions of $P$ and $Q$}

First consider two exact sequences of abelian groups
\begin{equation}\tag{3.8.1}
 0\longrightarrow L_1\xrightarrow{\ i\ }L_0\xrightarrow{\ u\ }P
 \longrightarrow0,\qquad
 0\longrightarrow M_1\xrightarrow{\ j\ }M_0\xrightarrow{\ v\ }Q
 \longrightarrow0.
\end{equation}
We shall describe the category of biextensions $E$ of $(P,Q)$ by an
abelian group $G$ for which the inverse image $(u,v)^*(E)$ is trivial;
that is, those for which there exists a homomorphism $s$ making the
following triangle commutative:
\begin{equation}\tag{3.8.2}
\begin{tikzcd}[column sep=huge,row sep=huge]
 E\arrow[d]\arrow[dr,dash,"s"]&\\
 P\times Q&L_0\times M_0\arrow[l,"u\times v"']
\end{tikzcd}
\end{equation}
and satisfying the additivity condition required in each of its two
arguments.  The choice of such an $s$ therefore identifies
$(u,v)^*(E)$ with the trivial biextension
$E_0=L_0\times M_0\times G$ of $(L_0,M_0)$ by $G$.  The left
exactness in $P,Q$ of the cofibred category of biextensions of
$(P,Q)$ by fixed $G$ (3.7.6) then shows that to give the biextension
$E$ of $(P,Q)$, together with the trivialization $s$ of its inverse
image, is essentially equivalent to giving two partial
trivializations of the trivial biextension $E_0$ of $(L_0,M_0)$ by
$G$, over $(L_1,M_0)$ and $(L_0,M_1)$ respectively, which agree over
$(L_1,M_1)$.  Since the biextensions induced by $E_0$ over
$(L_1,M_0)$, $(L_0,M_1)$, and $(L_1,M_1)$ identify with the trivial
biextensions, a partial trivialization over $(L_1,M_0)$ (respectively
over $(L_0,M_1)$) is equivalent to giving a homomorphism (or pairing)
\begin{equation}\tag{3.8.3}
 f:L_1\otimes M_0\longrightarrow G
 \qquad
 \bigl(\text{respectively }g:L_0\otimes M_1\longrightarrow G\bigr).
\end{equation}
The compatibility in question simply means that the two
homomorphisms agree over $L_1\otimes M_1$, namely
\begin{equation}\tag{3.8.4}
 f(x_1\otimes y_1)=g(x_1\otimes y_1),
\end{equation}
for sections $x_1,y_1$ of $L_1,M_1$ over an object of
$\underline T$.  In terms of the homomorphism $s$ of (3.8.2), the
forms $f$ and $g$ are characterized by
\begin{equation}\tag{3.8.5}
 s(x_1,y_0)=f(x_1\otimes y_0)e_{E/Q}(y_0),\qquad
 s(x_0,y_1)=g(x_0\otimes y_1)e_{E/P}(x_0),
\end{equation}
which again makes clear the necessity of condition (3.8.4).

Once the trivialization $s$ has been chosen, the other
trivializations are the morphisms
$s':L_0\times M_0\longrightarrow E$ of the form
\begin{equation}\tag{3.8.6}
 s'(x_0,y_0)=s(x_0,y_0)+h(x_0\otimes y_0),
\end{equation}
where
\[
 h:L_0\otimes M_0\longrightarrow G
\]
is an arbitrary form.  It is then clear from (3.8.5) that the forms
$f,g$ defined by $s$ are changed into
\begin{equation}\tag{3.8.7}
\begin{aligned}
 f'(x_1\otimes y_0)&=f(x_1\otimes y_0)+h(x_1\otimes y_0),\\
 g'(x_0\otimes y_1)&=g(x_0\otimes y_1)+h(x_0\otimes y_1).
\end{aligned}
\end{equation}
Thus we obtain:

\underline{Proposition} 3.8.8.
\underline{The subgroup of $\Biext^1(P,Q;G)$ formed by the classes
of biextensions whose inverse image under $(u,v)$ is a trivial
biextension of $(L_0,M_0)$} (that is, the kernel of
$\Biext^1(P,Q;G)\longrightarrow\Biext^1(L_0,M_0;G)$)
\underline{is canonically isomorphic, by the preceding description,
to the quotient of the group of pairs of forms $(f,g)$ in (3.8.3)
satisfying compatibility condition (3.8.4), by the pairs of forms
which arise by restriction from a form}
$h:L_0\otimes M_0\longrightarrow G$.

We now interpret this isomorphism in terms of homological algebra,
taking into account the canonical isomorphism
\[
 \Biext^1(P,Q;G)\simeq
 \Ext^1(P\mathbin{\overset{\mathbb L}{\otimes}}Q,G)
\]
of 3.6.5.  To this end, interpret the data (3.8.1) as two complexes
resolving $P$ and $Q$, respectively:
\begin{equation}\tag{3.8.9}
 L_\bullet:\ 0\longrightarrow L_1\longrightarrow L_0
 \longrightarrow0,\qquad
 M_\bullet:\ 0\longrightarrow M_1\longrightarrow M_0
 \longrightarrow0,
\end{equation}
and consider the tensor-product complex $L_\bullet\otimes M_\bullet$
\[
 0\longrightarrow L_1\otimes M_1
 \xrightarrow{d_2}
 L_1\otimes M_0+L_0\otimes M_1
 \xrightarrow{d_1}
 L_0\otimes M_0\longrightarrow0.
\]
There is, moreover, a canonical homomorphism in the derived category
\[
 P\mathbin{\overset{\mathbb L}{\otimes}}Q
 \longrightarrow L_\bullet\otimes M_\bullet,
\]
and hence an equally canonical homomorphism
\begin{equation}\tag{3.8.10}
 H^1\!\left(\Hom^{\cdot}(L_\bullet\otimes M_\bullet,G)\right)
 \longrightarrow
 \Ext^1(P\mathbin{\overset{\mathbb L}{\otimes}}Q,G).
\end{equation}
The group of $1$-cocycles of
$\Hom^{\cdot}(L_\bullet\otimes M_\bullet,G)$ identifies with the
group of pairs $(f,g)$ as in (3.8.3) which satisfy compatibility
condition (3.8.4): to such a pair one associates the homomorphism
$F:L_1\otimes M_0+L_0\otimes M_1\longrightarrow G$ with components
$(f,g)$,
\begin{equation}\tag{3.8.11}
 F(x_1\otimes y_0+x_0\otimes y_1)
 =f(x_1\otimes y_0)+g(x_0\otimes y_1).
\end{equation}
Moreover, the coboundary of the $0$-cochain
$h:L_0\otimes M_0\longrightarrow G$ is precisely the pair of its
restrictions to $L_1\otimes M_0$ and $L_0\otimes M_1$.  Consequently,
the group of pairs $(f,g)$ modulo the pairs $(h,h)$ described in
3.8.8 is canonically isomorphic to the first term of (3.8.10).
With this understood, one has the following compatibility.

\underline{Proposition} 3.8.12.
\underline{If}
\[
 \xi\in\Ker\!\left(
 \Biext^1(P,Q;G)\longrightarrow\Biext^1(L_0,M_0;G)\right)
\]
\underline{is described by the pair $(f,g)$ as in 3.8.8, then the
image of $\xi$ in}
$\Ext^1(P\mathbin{\overset{\mathbb L}{\otimes}}Q,G)$
\underline{under the isomorphism (3.6.5.1) is the image under
(3.8.10) of the element of its first term defined by the cocycle
(3.8.11)}.

The verification of this compatibility is left to the reader.

3.8.13. We shall apply this especially in the case
\[
 L_0=\mathbb Z[P],\qquad M_0=\mathbb Z[Q],
\]
where $L_1$ and $M_1$ are respectively the kernels of the canonical
homomorphisms $\mathbb Z[P]\longrightarrow P$ and
$\mathbb Z[Q]\longrightarrow Q$.  Notice the canonical isomorphisms
\[
 L_0\mathbin{\overset{\mathbb L}{\otimes}}M_0
 \simeq L_0\otimes M_0
 \simeq\mathbb Z[P\times Q].
\]
Since $L_0$ and $M_0$ are flat, (1.4.5) gives canonical isomorphisms
\[
 \Ext^i(L_0\mathbin{\overset{\mathbb L}{\otimes}}M_0,G)
 \simeq H^i(P\times Q,G_{P\times Q}).
\]
It follows from this and 3.6.5, or directly from the definitions and
1.4, that the functor which assigns to each biextension of
$(L_0,M_0)$ by $G$ the $G_{P\times Q}$-torsor over $P\times Q$ that
it induces is an equivalence of categories.  (This remains valid
without using an additive structure on $P,Q$.)  Consequently, if
$E$ is a biextension of $(P,Q)$ by $G$, giving a trivialization of
its inverse image over $(L_0,M_0)$ is equivalent to giving a section
of $E$ considered only as an object over $P\times Q$, that is, a
trivialization of $E$ considered as a torsor under $G_{P\times Q}$.

Biextensions equipped with this additional structure are therefore
described by the pairs $(f,g)$ of (3.8.3) satisfying compatibility
(3.8.4), and their classification up to isomorphism is the one given
in 3.8.8.

With the notation of 3.5, one may also regard $L_1$ as the cokernel
of $L_2(P)\longrightarrow L_1(P)$, and similarly for $M_1$.  This
makes it possible to express the data $(f,g)$ in terms of a system
of homomorphisms
\[
\left\{
\begin{aligned}
 f'&:L_1(P)\otimes L_0(Q)
     =\mathbb Z(P\times P\times Q)\longrightarrow G,\\
 g'&:L_0(P)\otimes L_1(Q)
     =\mathbb Z(P\times Q\times Q)\longrightarrow G,
\end{aligned}
\right.
\]
that is, homomorphisms of sheaves of \underline{sets}
\[
 P\times P\times Q\longrightarrow G,\qquad
 P\times Q\times Q\longrightarrow G,
\]
satisfying five compatibility conditions.  These express that $f'$
and $g'$ factor through $f$ and $g$---which gives $2\times2$
equations, in view of the expression of $L_2(P)$ and $L_2(Q)$ as
sums of two terms---together with condition (3.8.4).  One then sees
that $f'$ and $g'$ are precisely the partial multiplication data
defining a biextension, and that the five conditions in question are
exactly the five axioms for biextensions considered in 2.1 (cf. the
proof of 3.6.4).

\subsection*{Bibliography to Expos\'e VII}

\begin{enumerate}[label={[\arabic*]}]
\item M. Demazure, J. Giraud, and M. Raynaud,
\emph{Schémas abéliens}, Séminaire de la Faculté des Sciences
d'Orsay 1967/68 (to appear).
\item J. Giraud, \emph{Cohomologie non abélienne}, thesis, Paris,
1967 (to appear).
\item A. Grothendieck, \emph{Sur quelques points d'Algèbre
homologique}, Tohoku Mathematical Journal 9 (1956), 119--221.
\item A. Grothendieck, \emph{Catégories cofibrées additives et
Complexe cotangent relatif}, Lecture Notes in Mathematics 79,
Springer, 1968.
\item J.-L. Verdier, \emph{Catégories dérivées des catégories
abéliennes} (to appear).
\item D. Mumford, \emph{Biextensions of formal groups}, Tata
Institute of Fundamental Research, 1968 (to appear).
\end{enumerate}
